% 原书第六篇6.6节249--258，印刷260--261，原PDF269--270；
% 书末答案349--350，PDF358--359。
\nuclearreactionproblem{249}{260--261}
动能 $T_\alpha=7.0\,\mathrm{MeV}$ 的
$\alpha$ 粒子与起初静止的 $\mathrm{^{6}Li}$ 核
弹性散射。原题称“两粒子出射方向夹角”
$\theta=60^\circ$，求反冲锂核动能。
\exsolution
令入射粒子质量 $m=4u$、靶核质量 $M=6u$，
散射后两个粒子动量模为 $p'_a,p_R$，
它们之间的实际夹角为 $\gamma$。
动量、能量守恒分别为
\[
 p_0^2=p_a'^2+p_R^2+2p_a'p_R\cos\gamma,
 \qquad
 \frac{p_0^2}{2m}=
 \frac{p_a'^2}{2m}+\frac{p_R^2}{2M}.
\]
两式消去 $p_0^2$：
\[
 2p_a'p_R\cos\gamma=
 \left(\frac mM-1\right)p_R^2.
\]
因 $m<M$，实际 $\cos\gamma<0$，
所以两出射动量夹角必须\emph{大于}$90^\circ$。
若原题角度 $60^\circ$ 是实际出射夹角，
则不存在物理解。
原书答案 $T_R=6.0\,\mathrm{MeV}$ 可由补角
$\gamma=120^\circ$ 求出：
利用
$p_a'^2=2m(T_\alpha-T_R)$、
$p_R^2=2MT_R$，得
\[
 \boxed{T_R=
 \frac{T_\alpha}{1+
 \frac{(M-m)^2}{4mM\cos^2\gamma}}=
 6.0\,\mathrm{MeV}}.
\]
\excheck \emph{印刷60°与解6.0MeV不同时满足带符号动量守恒}，
因为6MeV实际要求夹角120°。
本稿同时列出条件及书末结果，不悄悄修改出射角。

\nuclearreactionproblem{250}{261}
中子与起初静止的氘核弹性碰撞。
求中子损失动能占初始动能的比例：
(a)正面碰撞；(b)中子偏转 $90^\circ$。
\exsolution
令中子质量 $m$、氘核质量 $M\simeq2m$。
(a) 质量为$m$的粒子撞静止质量$M$物体，
弹性正碰后动能比例
$T'_n/T_n=[(m-M)/(m+M)]^2$，
所以
\[
 \boxed{\eta_a=
 1-\left(\frac{m-M}{m+M}\right)^2
 =\frac{4mM}{(m+M)^2}=\frac89}.
\]
(b) 末态中子与入射方向正交，
矢量动量守恒给
$p_R^2=p_0^2+p_n'^2$。
代能量守恒
$p_0^2/(2m)=p_n'^2/(2m)+p_R^2/(2M)$，
得到
$p_n'^2/p_0^2=(M-m)/(M+m)=1/3$，故
\[
 \boxed{\eta_b=1-\frac13=\frac23}.
\]
\excheck (b) $90^\circ$ 指中子偏转角，
不是中子和氘核两个出射粒子的夹角。

\nuclearreactionproblem{251}{261}
氘核与起初静止的质子弹性碰撞，
求氘核在实验室系中可能达到的最大散射角。
\exsolution
入射氘核质量$m_d\simeq2m_p$，
靶质子质量$M=m_p$。
质心系中入射氘核速度
$u^*=Mv_0/(m_d+M)$，
质心平移速度
$V=m_dv_0/(m_d+M)$。
实验室末态氘核速度矢量是
以 $V$ 为圆心、半径 $u^*$ 的速度圆。
从原点向该圆作切线时实验室偏转角最大：
\[
 \boxed{\sin\theta_{\max}=\frac{u^*}{V}
 =\frac M{m_d}\simeq\frac12},
 \quad
 \boxed{\theta_{\max}=30^\circ}.
\]
\excheck 重入射粒子撞轻靶核时，
其实验室偏角有限制，并非可取任意后散射角。

\nuclearreactionproblem{252}{261}
原书将核半径写为
$R=0.13A^{1/3}\,\mathrm{pm}$，
$A$ 为质量数。按此题面估算原子核质量密度
和核子数密度，并核对书末量级。
\exsolution
若严格采用原印
$r_0=0.13\,\mathrm{pm}=1.3\times10^{-13}\,\mathrm m$，
半径为$r_0A^{1/3}$，则
\[
 \rho=\frac{Au}{(4\pi/3)r_0^3A}
 =\frac{3u}{4\pi r_0^3},\quad
 n=\frac{3}{4\pi r_0^3}.
\]
数值约为
\[
 \boxed{\rho_{\rm literal}\simeq
 1.8\times10^{5}\,\mathrm{kg/cm^3}},\qquad
 \boxed{n_{\rm literal}\simeq
 1.1\times10^{32}\,\mathrm{nucleons/cm^3}}.
\]
这与书末给出的约
$2\times10^{11}\,\mathrm{kg/cm^3}$、
$10^{38}\,\mathrm{nucleons/cm^3}$
相差$10^6$。
若采用物理上公认的核尺度
$r_0=1.3\,\mathrm{fm}=0.0013\,\mathrm{pm}$，
则
\[
 \boxed{\rho_{\rm nuclear}\simeq
 2\times10^{11}\,\mathrm{kg/cm^3}},\quad
 \boxed{n_{\rm nuclear}\simeq
 10^{38}\,\mathrm{nucleons/cm^3}},
\]
恢复原书答案量级。
\excheck 原图确印$0.13A^{1/3}\,\mathrm{pm}$，
而从书末密度可判断物理上需要
$0.0013A^{1/3}\,\mathrm{pm}$；
本译稿\emph{并列两种数据}并列入勘误。

\nuclearreactionproblem{253}{261}
利用质量数和电荷数守恒补全：
(a) $\mathrm{^{10}B}(x,\alpha)\mathrm{^{8}Be}$；
(b) $\mathrm{^{17}O}(d,n)x$；
(c) $\mathrm{^{23}Na}(p,x)\mathrm{^{20}Ne}$；
(d) $x(p,n)\mathrm{^{37}Ar}$。
\exsolution
核反应的核子总数$A$和电荷总数$Z$守恒。
(a) 输入粒子应有$A=8+4-10=2$、
$Z=4+2-5=1$，即
$\boxed{x=d=\mathrm{^{2}H}}$。
(b) 输出
$A=17+2-1=18$、
$Z=8+1-0=9$，即
$\boxed{x=\mathrm{^{18}F}}$。
(c) 输入后$A=24,Z=12$，
$\mathrm{^{20}Ne}$有$A=20,Z=10$，
剩余 $\boxed{x=\alpha=\mathrm{^{4}He}}$。
(d) 输入靶核$A=37+1-1=37$、
$Z=18-1=17$，
故$\boxed{x=\mathrm{^{37}Cl}}$。
\excheck 原书简答(b)纯文本出现
$\mathrm{^{17}F}$，但题面明确是
$\mathrm{^{17}O}(d,n)$，
按质量数守恒产物必须是$\mathrm{^{18}F}$；
已记录底本可能的上标排印或识别差异。

\nuclearreactionproblem{254}{261}
证明核的结合能可按质量超额计算：
$E_b=[Z\Delta_H+(A-Z)\Delta_n-\Delta_{\rm at}]c^2$。
\exsolution
若用中性原子的实测质量$M_{\rm at}$，
核结合能（忽略电子化学结合能）为
\[
 E_b=[Zm_H+(A-Z)m_n-M_{\rm at}]c^2.
\]
把原子质量与中子质量写成
$m_H=u+\Delta_H$、
$m_n=u+\Delta_n$、
$M_{\rm at}=Au+\Delta_{\rm at}$。
则
\[
 Zm_H+(A-Z)m_n-M_{\rm at}
 =Z\Delta_H+(A-Z)\Delta_n-\Delta_{\rm at},
\]
其中$Zu+(A-Z)u-Au=0$，因此
\[
 \boxed{E_b=
 [Z\Delta_H+(A-Z)\Delta_n-\Delta_{\rm at}]c^2}.
\]
\excheck 选用\emph{氢原子}质量而非裸质子质量，
正是为了抵消相应的原子电子静质量。

\nuclearreactionproblem{255}{261}
某核的质子数与中子数相等，
且半径比 $\mathrm{^{27}Al}$ 核小1.5倍。
求核种与其结合能。
\exsolution
核半径比例$R\propto A^{1/3}$，
所以
\[
 \frac{R_x}{R_{\mathrm{Al}}}=\frac1{1.5}
 \Longrightarrow A_x=\frac{27}{1.5^3}=8.
\]
质子数等于中子数，故$Z=8/2=4$，
该核为$\boxed{\mathrm{^{8}Be}}$。
按原书附录质量超额
$\Delta_H=0.00783u$、
$\Delta_n=0.00867u$、
$\Delta_{\mathrm{^{8}Be}}=0.00531u$，
\[
 \boxed{E_b=
 [4(0.00783)+4(0.00867)-0.00531]
 (931.5\,\mathrm{MeV})
 \simeq56.5\,\mathrm{MeV}}.
\]
\excheck $\mathrm{^{8}Be}$ 对裂变为两$\alpha$
粒子是不稳定的，但仍可计算它的核结合能。

\nuclearreactionproblem{256}{261}
用原子质量表求：
(a) $\mathrm{^{16}O}$ 平均每核子结合能；
(b) $\mathrm{^{11}B}$ 中单中子及 $\alpha$
的分离能；(c) 把 $\mathrm{^{16}O}$ 分成四个
相同粒子所需能量。
\exsolution
原书质量超额（单位$u$）：
$\Delta_H=0.00783$、
$\Delta_n=0.00867$、
$\Delta_{\mathrm{^{16}O}}=-0.00509$、
$\Delta_{\mathrm{^{11}B}}=0.00930$、
$\Delta_{\mathrm{^{10}B}}=0.01294$、
$\Delta_{\mathrm{^{7}Li}}=0.01601$、
$\Delta_\alpha=0.00260$。
(a)
\[
 \boxed{\frac{E_b(\mathrm{^{16}O})}{16}
 =\frac{[8\Delta_H+8\Delta_n-\Delta_O]c^2}{16}
 \simeq8.0\,\mathrm{MeV}}.
\]
(b) 分离中子反应
$\mathrm{^{11}B}\to\mathrm{^{10}B}+n$，
以及分离$\alpha$反应
$\mathrm{^{11}B}\to\mathrm{^{7}Li}+\alpha$，故
\[
 \boxed{S_n=(\Delta_{\mathrm{^{10}B}}
 +\Delta_n-\Delta_{\mathrm{^{11}B}})c^2
 \simeq11.5\,\mathrm{MeV}},
\]
\[
 \boxed{S_\alpha=(\Delta_{\mathrm{^{7}Li}}
 +\Delta_\alpha-\Delta_{\mathrm{^{11}B}})c^2
 \simeq8.7\,\mathrm{MeV}}.
\]
(c) 四个相同粒子应是$\alpha$，
\[
 \boxed{E= [4\Delta_\alpha-\Delta_O]c^2
 \simeq14.4\,\mathrm{MeV}\approx14.5\,\mathrm{MeV}}.
\]
\excheck 分离能是末态总静能减母核静能；
守恒的核子总数使整数原子质量数自动抵消。

\nuclearreactionproblem{257}{261}
求 $\mathrm{^{11}B}$ 核中中子与质子的
分离能之差，并解释其原因。
\exsolution
中子分离能
$S_n=[M(\mathrm{^{10}B})+m_n-M(\mathrm{^{11}B})]c^2$。
质子分离能
$S_p=[M(\mathrm{^{10}Be})+m_H-M(\mathrm{^{11}B})]c^2$
（采用中性原子质量）。
两式相减
\[
 S_n-S_p=
 [M(\mathrm{^{10}B})+m_n
 -M(\mathrm{^{10}Be})-m_H]c^2.
\]
代原书质量超额：
$\Delta_{\mathrm{^{10}B}}=0.01294$、
$\Delta_{\mathrm{^{10}Be}}=0.01354$、
$\Delta_n=0.00867$、
$\Delta_H=0.00783$，
\[
 \boxed{S_n-S_p=
 (0.01294+0.00867-0.01354-0.00783)
 (931.5\,\mathrm{MeV})
 \simeq0.22\,\mathrm{MeV}}.
\]
\excheck 差异来自质子受库仑排斥的附加效应，
同时受娘核结构差异影响，不能仅归因于核子静质量。

\nuclearreactionproblem{258}{261}
$\mathrm{^{20}Ne}$、
$\alpha$ 与 $\mathrm{^{12}C}$
平均每核子结合能分别为
$8.03$、$7.07$、$7.68\,\mathrm{MeV}$。
求将$\mathrm{^{20}Ne}$分为
两个$\alpha$与一个$\mathrm{^{12}C}$所需能量。
\exsolution
初核的结合能
$E_i=20(8.03)\,\mathrm{MeV}$。
末态三个核的总结合能
$E_f=2\cdot4(7.07)+12(7.68)$ MeV。
分裂所需能量为失去的结合能差
\[
 \boxed{E_{\rm req}=E_i-E_f
 =20(8.03)-8(7.07)-12(7.68)
 \simeq11.9\,\mathrm{MeV}}.
\]
\excheck 结合能越大质量亏损越大；
若初态结合能大于产物和，
分裂必须吸收正能量。
